\documentclass{../notatki}
\fancyhead[C]{}

\title{Problemy}

\begin{document}

\section{Zestaw 1}

\subsection{Zadanie 1}

\subsubsection{Udowodnij, że każda skończona macierz kwadratowa $A$ nad ciałem
  liczb zespolonych jest podobna do swojej transpozycji
\texorpdfstring{$A^T$}{A\^T}}

Każda macierz $A$ jest podobna do swojej postaci Jordana $A \sim J_A$.
Ponieważ podobieństwo jest tranzytywne:
$$
A \sim B, B \sim C \Rightarrow B = P_A A P_A^{-1}, C = P_B B P_B^{-1}
\Rightarrow C =  P_B P_A A P_A^{-1} P_B^{-1} \Rightarrow C = P_C A P_C^{-1}
\rightarrow A \sim C
$$
to jeśli $J_A \sim J_{A^T}$, to $A \sim A^T$, ponieważ:
$$
A \sim J_A \sim J_{A^T} \sim A^T
$$

Macierz Jordana $J_A$ składa się z bloków, dla odpowiedniej wartości
własnej $\lambda$:
$$
J_k(\lambda) =
\begin{bmatrix}
  \lambda & 1 & 0 & \hdots & 0 \\
  0 & \lambda & 1 & \hdots & 0 \\
  \vdots & & \ddots & \ddots & \vdots \\
  0 & 0 & 0 & 0 & \lambda
\end{bmatrix}
\quad J_A = \textrm{diag}(J_1(\lambda_1), \dots, J_k(\lambda_k))
\quad J_k^T(\lambda) =
\begin{bmatrix}
  \lambda & 0 & \hdots & 0 \\
  1 & \lambda & \hdots & 0 \\
  0 & 1 & \lambda & 0 \\
  \vdots & \vdots& \ddots & \vdots \\
  0 & 0 & 0 & \lambda
\end{bmatrix}
$$
Z własności złożenia $\textrm{diag}$, wiemy, że $J_A^T =
\textrm{diag}(J_1^T(\lambda_1), \dots, J_k^T(\lambda_k))$

Macierz antydiagonalna $R$, ma dwie kluczowe własności. Po pierwsze,
transformuje nadprzekątną macierzy, na podprzekątną, jak na poniższym
przykładzie:
$$
\begin{bmatrix}
  0 & 1 \\
  1 & 0
\end{bmatrix}
\begin{bmatrix}
  \alpha & \beta \\
  \gamma & \delta
\end{bmatrix}
\begin{bmatrix}
  0 & 1 \\
  1 & 0
\end{bmatrix}
=
\begin{bmatrix}
  \gamma & \delta \\
  \alpha & \beta
\end{bmatrix}
\begin{bmatrix}
  0 & 1 \\
  1 & 0
\end{bmatrix}
=
\begin{bmatrix}
  \delta & \gamma \\
  \beta & \alpha
\end{bmatrix}
$$
$R$ odwraca kolejność wierszy i kolumn, więc element na pozycji
$(i,i+1)$ przechodzi na $(k-i+1, k-(i+1)+1)=(k-i+1, k-i)$, czyli z
nadprzekątnej na podprzekątną.
Po drugie; $R^{-1} = R$, zatem:
$$
RJ_k(\lambda)R^{-1} = RJ_k(\lambda)R = J_k^T(\lambda)
$$
stąd wynika, że $J_k(\lambda) \sim J_k^T(\lambda)$, ponieważ macierz $R$
spełnia formalny warunek podobieństwa. Jeśli skonstruujemy macierz $Q$,
taką, że $Q = \textrm{diag}(R_1, \dots, R_k)$, to:
$$
Q^{-1} J_A Q = J_A^T \sim J_A \sim J_{A^T} \sim A^T \sim A \quad \square
$$

\subsubsection{Wyznacz liczbę klas podobieństwa macierzy 2×2 nad ciałem dwu-
elementowym \texorpdfstring{$\mathbb{F}_2$}{F\_2}}

Niech $A \in F_2^{2 \times 2}$. Oznaczmy przez
$$
p_A(t) = \det(tI-A)
$$
wielomian charakterystyczny, a przez $m_A(t)$ wielomian minimalny macierzy $A$.
Z twierdzenia Cayleya--Hamiltona mamy $p_A(A)=0$, więc
$$
m_A(t)\mid p_A(t).
$$
Dla macierzy $2\times 2$ zachodzi ponadto: klasy podobieństwa są wyznaczone
przez (równoważnie) racjonalną postać kanoniczną, czyli przez możliwe pary
$(p_A,m_A)$ (inaczej: przez wielomiany niezmiennicze).

Istnieją dokładnie cztery możliwości dla $p_A(t)$:
$$
t^2,\qquad t^2+1=(t+1)^2,\qquad t^2+t=t(t+1),\qquad t^2+t+1.
$$
Rozpatrzmy je kolejno.

\noindent\textbf{(1) $p_A(t)=t^2$.}
Wtedy jedyną wartością własną (w $\overline{\mathbb{F}_2}$) jest $0$ i
$m_A(t)\mid t^2$, więc $m_A(t)\in\{t,t^2\}$.
\begin{itemize}
  \item Jeśli $m_A(t)=t$, to $A=0$ (jedna klasa podobieństwa).
  \item Jeśli $m_A(t)=t^2$, to $A$ jest niezerowa nilpotentna i jest podobna do
    $$
    N =
    \begin{pmatrix}
      0 & 1 \\
      0 & 0
    \end{pmatrix}
    $$
\end{itemize}
W tym przypadku dostajemy $2$ klasy.

\noindent\textbf{(2) $p_A(t)=t^2+1=(t+1)^2$.}
Wtedy $m_A(t)\mid (t+1)^2$, więc $m_A(t)\in\{t+1,(t+1)^2\}$.
\begin{itemize}
  \item Jeśli $m_A(t)=t+1$, to $A=I$ (jedna klasa).
  \item Jeśli $m_A(t)=(t+1)^2$, to $A$ ma nietrywialny blok Jordana dla
    $\lambda=1$ i jest podobna do
    $$
    J =
    \begin{pmatrix}
      1 & 1 \\
      0 & 1
    \end{pmatrix}
    $$
    (druga klasa).
\end{itemize}
W tym przypadku dostajemy $2$ klasy.

\noindent\textbf{(3) $p_A(t)=t^2+t=t(t+1)$.}
Wielomian charakterystyczny rozkłada się na iloczyn dwóch różnych czynników
liniowych, więc $A$ ma dwie różne wartości własne $0$ i $1$ w $\mathbb{F}_2$,
a zatem jest diagonalizowalna i podobna do
$$
D =
\begin{pmatrix}
  0 & 0 \\
  0 & 1
\end{pmatrix}.
$$
Daje to dokładnie $1$ klasę.

\noindent\textbf{(4) $p_A(t)=t^2+t+1$.}
Wielomian $t^2+t+1$ jest nierozkładalny nad $\mathbb{F}_2$, więc jedyny
możliwy wielomian minimalny to $m_A(t)=\chi_A(t)$.
Wtedy $A$ jest podobna do macierzy towarzyszącej (kompanionowej) tego
wielomianu:
$$
C =
\begin{pmatrix}
  0 & 1 \\
  1 & 1
\end{pmatrix}.
$$
Daje to dokładnie $1$ klasę.

Sumując liczbę klas we wszystkich przypadkach, otrzymujemy
$$
2+2+1+1=6.
$$
Zatem liczba klas podobieństwa macierzy $2\times 2$ nad $\mathbb{F}_2$ wynosi 6.

\subsection{Udowodnij twierdzenie Gerszgorina}

\textit{Każda wartość własna $\lambda$ macierzy $A \in \mathbb{C}^{n \times n}$
  leży wewnątrz lub na brzegu jednego z kół o środku $a_{ii}$ i
  promieniu $R_i = \sum_{j \neq i} |a_{ij}|$
}

Dla $A \ket{\lambda} = \lambda \ket{\lambda}$:
$$
\sum_{j = 1}^n a_{ij} \ket{\lambda}_j = \lambda \ket{\lambda}_i
$$
$$
a_{ii}\ket{\lambda}_i + \sum_{j \neq i} a_{ij} \ket{\lambda}_j =
\lambda \ket{\lambda}_i
\Rightarrow (\lambda - a_{ii})\ket{\lambda}_i = \sum_{j \neq i}
a_{ij} \ket{\lambda}_j
$$
Jeśli $\ket{\lambda}_i \neq 0$, to dzielimy przez $\ket{\lambda}_i$ i
bierzemy moduł:
$$
|\lambda - a_{ii}| = \left|\sum_{j \neq i} a_{ij}
\frac{\ket{\lambda}_j}{\ket{\lambda}_i}\right| \leq \sum_{j \neq i}
|a_{ij}| \left|\frac{\ket{\lambda}_j}{\ket{\lambda}_i}\right| =
\left|\frac{1}{\ket{\lambda}_i}\right|\sum_{j \neq i} \left|a_{ij}
\ket{\lambda}_j\right|
$$

Istnieje takie $k$, że:
$$
\left|\ket{\lambda}_k\right| = \max_{1 \leq i \leq n}
\left|\ket{\lambda}_i\right|
$$
Dla $i = k$:
$$
\left|\lambda - a_{kk}\right| \leq
\left|\frac{1}{\ket{\lambda}_k}\right|\sum_{j \neq k} \left|a_{kj}
\ket{\lambda}_j\right|
$$
Z maksymalności $\left|\ket{\lambda}_k\right|$ mamy
$\left|\ket{\lambda}_j\right|
\leq \left|\ket{\lambda}_k\right|$, więc $\left|a_{kj}
\ket{\lambda}_j\right| \leq |a_{kj}|
\left|\ket{\lambda}_k\right|$, co podstawiając:
$$
\left|\lambda - a_{kk}\right| \leq \sum_{j \neq k} |a_{kj}| = R_k
$$
Jest to typowa konstrukcja koła Gerszgorina. $\square$

\subsection{Zadanie 3}

Dla $T_1: V \to W$, oraz $T_2: W \to V$, będą dwoma przekształceniami liniowymi.
$$
T = T_1 \circ T_2: V \to V
$$
Prawdziwym jest, że $\ker T_1 = \ker T$, oraz $\ker T_2 \subset \Im T_1$,
czy $T_2$ jest różnowartościowe?

$T_2$ jest różnowartościowe, jeśli $\ker T_2 = \{0\}$. Zatem, chcemy udowodnić,
że $\forall_{w \in \ker T_2} w = 0$. Wiemy, że $w \in \Im T_1$. Istnieje więc,
takie $v \in V : T_1(v) = w$. Z definicji $w$, mamy:
$$
T(v) = T_2(T_1(v)) = T_2(w) = 0
$$
więc $v \in \ker T = \ker T_1$. Zatem:
$$
w = T_1(v) = 0.
$$
Dowolny wektor $v \in V$, po przekształceniu $T_1(v)$ daje nam
dowolny $w \in W$, który musi być zerowy. $\square$

\section{Zestaw 2}

\textit{
  Stanem absolutnie maksymalnie splątanym (AME state) nazywamy taki
  stan skwantowy, którego wszystkie zredukowane macierze gęstości są
  maksymalnie mieszane.
}

\subsection{Zadanie 1}

\subsubsection{GHZ}

$$
\ket{GHZ} = \frac{1}{\sqrt{2}} \left(\ket{000} + \ket{111}\right)
$$
$$
\rho = \ket{GHZ} \bra{GHZ} = \frac{1}{2} \left(\ket{000} \bra{000} +
\ket{000} \bra{111} + \ket{111} \bra{000} + \ket{111} \bra{111}\right)
$$
$$
\Tr_{BC}(\ket{000}\bra{000}) = \ket{0}\bra{0} \quad
\Tr_{BC}(\ket{111}\bra{111}) = \ket{1}\bra{1}
$$
$$
\Tr_{BC}(\ket{000}\bra{111}) = \Tr_{BC}\Big( (\ket{0}\bra{1})_A
\otimes (\ket{00}\bra{11})_{BC} \Big)
=
(\ket{0}\bra{1})\,\Tr(\ket{00}\bra{11})
=
\ket{0}\bra{1}\,\braket{11}{00}
=
\ket{0}\bra{1}\cdot 0
=
0
$$
$$
\Tr_{BC}(\ket{111}\bra{000})
=
\ket{1}\bra{0}\,\braket{00}{11}
=
\ket{1}\bra{0}\cdot 0
=
0
$$

$$
\rho_A = \frac{1}{2}\left(1, 0, 0, 1\right) = \rho_B = \rho_C
$$

Każda jedno-kubitowa redukcja jest maksymalnie mieszana, więc
$\ket{GHZ}$ jest absolutely maximally entangled.

\subsubsection{W}

Czy stan $\ket{W}$ jest absolutnie maksymalnie splątany?
$$
\ket{W} = \frac{1}{\sqrt{3}}\left(\ket{001} + \ket{010} + \ket{100}\right)
$$
$$
\rho = \ket{W} \bra{W}
$$

$$
\rho_A = \frac{1}{3} \left(\ket{0}\bra{0} \braket{01}{01} +
  \ket{0}\bra{0} \braket{10}{10} + \ket{1}\bra{1}
\braket{00}{00}\right) = \frac{2}{3} \ket{0}\bra{0} + \frac{1}{3} \ket{1}\bra{1}
$$

Z symetrii stanu $\rho_A = \rho_B = \rho_C$. Ponieważ redukcja do jednego kubitu
nie jest maksymalnie mieszana, stan $\ket{W}$ nie jest absolutnie
maksymalnie splątany.

\subsection{Zadanie 2}

$$
\ket{\psi} = \frac{1}{\sqrt{2}} \left(\ket{0} + i\ket{1}\right)
$$

\subsubsection{Podaj współrzędne stanu}

$$
\ket{\psi} = \cos\left(\frac{\theta}{2}\right)\ket{0} +
e^{i\phi}\sin\left(\frac{\theta}{2}\right)\ket{1}
$$
$$
\cos\left(\frac{\theta}{2}\right) = \frac{1}{\sqrt{2}} \quad
e^{i\phi} \sin\left(\frac{\theta}{2}\right) = \frac{i}{\sqrt{2}}
$$
$$
\theta = \frac{\pi}{2} \quad \phi = \frac{\pi}{2}
$$

Wektor w sferze blocha ma współrzędne:
$$
(x, y, z) = (\sin\theta\cos\phi, \sin\theta\sin\phi, \cos\theta) = (0, 1, 0)
$$
$$
\langle X \rangle = 0 \quad \langle Y \rangle = 1 \quad \langle Z \rangle = 0
$$

\subsubsection{Pomiar uogólniony}

$$
M_i = \frac{\ket{\psi_i}\bra{\psi_i}}{2}
$$
$$
M_1 = \frac{\ket{0}\bra{0}}{2} = \frac{1}{2}
\begin{bmatrix}
  1 \\ 0
\end{bmatrix}
\begin{bmatrix}
  1 & 0
\end{bmatrix} =
\frac{1}{2}
\begin{bmatrix}
  1 & 0
  \\ 0 & 0
\end{bmatrix}
$$
$$
M_2 = \frac{
  \left(\frac{1}{\sqrt{3}} \ket{0} + \frac{\sqrt{2}}{\sqrt{3}} \ket{1}\right)
  \left(\frac{1}{\sqrt{3}} \bra{0} + \frac{\sqrt{2}}{\sqrt{3}} \bra{1}\right)
}{2} = \frac{1}{6}
\begin{bmatrix}
  1 & \sqrt{2} \\
  \sqrt{2} & 2
\end{bmatrix}
$$
$$
M_3 = \frac{
  \left(\frac{1}{\sqrt{3}} \ket{0} + \frac{\sqrt{2}}{\sqrt{3}}
  e^{\frac{2\pi}{3}} \ket{1}\right)
  \left(\frac{1}{\sqrt{3}} \bra{0} + \frac{\sqrt{2}}{\sqrt{3}}
  e^{\frac{2\pi}{3}} \bra{1}\right)
}{2} = \frac{1}{6}
\begin{bmatrix}
  1 & \sqrt{2}e^{\frac{2\pi}{3}} \\
  \sqrt{2}e^{\frac{2\pi}{3}} & 2e^{\frac{4\pi}{3}}
\end{bmatrix}
$$
$$
M_4 = \frac{1}{6}
\begin{bmatrix}
  1 & \sqrt{2}e^{\frac{4\pi}{3}} \\
  \sqrt{2}e^{\frac{4\pi}{3}} & 2e^{\frac{8\pi}{3}}
\end{bmatrix}
$$

$$
p_i = \bra{\psi} M_i \ket{\psi} \quad \ket{\psi} = \frac{1}{\sqrt{2}}
\begin{bmatrix}
  1 \\ -i
\end{bmatrix}
\quad \bra{\psi} = \frac{1}{\sqrt{2}}
\begin{bmatrix}
  1 & -i
\end{bmatrix}
$$
$$
p_1 = \frac{1}{4}
\begin{bmatrix}
  1 & -1
\end{bmatrix}
\begin{bmatrix}
  1 & 0
  \\ 0 & 0
\end{bmatrix}
\begin{bmatrix}
  1 \\ -i
\end{bmatrix} = \frac{1}{4}
\begin{bmatrix}
  1 & 0
\end{bmatrix}
\begin{bmatrix}
  1 \\ -i
\end{bmatrix} = \frac{1}{4}
$$
$$
p_2 = \frac{1}{12}
\begin{bmatrix}
  1 & -1
\end{bmatrix}
\begin{bmatrix}
  1 & \sqrt{2} \\
  \sqrt{2} & 2
\end{bmatrix}
\begin{bmatrix}
  1 \\ -i
\end{bmatrix} = \frac{1}{12}
\begin{bmatrix}
  1 - \sqrt{2} & \sqrt{2} - 2
\end{bmatrix}
\begin{bmatrix}
  1 \\ -i
\end{bmatrix} = \frac{1 - \sqrt{2} -i(\sqrt{2} - 2)}{12}
$$
$$
p_3 = \frac{1}{12}
\begin{bmatrix}
  1 & -1
\end{bmatrix}
\begin{bmatrix}
  1 & \sqrt{2}e^{\frac{2\pi}{3}} \\
  \sqrt{2}e^{\frac{2\pi}{3}} & 2e^{\frac{4\pi}{3}}
\end{bmatrix}
\begin{bmatrix}
  1 \\ -i
\end{bmatrix} = \frac{1}{12}
\begin{bmatrix}
  1 - \sqrt{2}e^{\frac{2\pi}{3}} & \sqrt{2}e^{\frac{2\pi}{3}} -
  2e^{\frac{4\pi}{3}}
\end{bmatrix}
\begin{bmatrix}
  1 \\ -i
\end{bmatrix} = \frac{1 - \sqrt{2}e^{\frac{2\pi}{3}} -
i(\sqrt{2}e^{\frac{2\pi}{3}} - 2e^{\frac{4\pi}{3}})}{12}
$$
$$
\frac{1}{12}
\begin{bmatrix}
  1 & -1
\end{bmatrix}
\begin{bmatrix}
  1 & \sqrt{2}e^{\frac{4\pi}{3}} \\
  \sqrt{2}e^{\frac{4\pi}{3}} & 2e^{\frac{8\pi}{3}}
\end{bmatrix}
\begin{bmatrix}
  1 \\ -i
\end{bmatrix} = \frac{1 - \sqrt{2}e^{\frac{4\pi}{3}} -
i(\sqrt{2}e^{\frac{4\pi}{3}} - 2e^{\frac{8\pi}{3}})}{12}
$$

\subsection{Zadanie 3}

$$
C, S : \mathcal{H}_A\otimes\mathcal{H}_B \to \mathcal{H}_A\otimes\mathcal{H}_B
$$
$$
C =
\begin{bmatrix}
  1 & 0 & 0 & 0 \\
  0 & 1 & 0 & 0 \\
  0 & 0 & 0 & 1 \\
  0 & 0 & 1 & 0
\end{bmatrix} \quad
S =
\begin{bmatrix}
  1 & 0 & 0 & 0 \\
  0 & 0 & 1 & 0 \\
  0 & 1 & 0 & 0 \\
  0 & 0 & 0 & 1
\end{bmatrix}
$$

\subsubsection{a}

$$
H_1 \otimes H_2 = \frac{1}{2}
\begin{bmatrix}
  1 \cdot H & 1 \cdot H \\
  1 \cdot H & -1 \cdot H
\end{bmatrix} = \frac{1}{2}
\begin{bmatrix}
  1 & 1 & 1 & 1 \\
  1 & -1 & 1 & -1 \\
  1 & 1 & -1 & -1 \\
  1 & -1 & -1 & 1
\end{bmatrix}
$$
$$
(H_1 \otimes H_2)C_{12}(H_1 \otimes H_2) = \frac{1}{4}
\begin{bmatrix}
  1 & 1 & 1 & 1 \\
  1 & -1 & 1 & -1 \\
  1 & 1 & -1 & -1 \\
  1 & -1 & -1 & 1
\end{bmatrix}
\begin{bmatrix}
  1 & 0 & 0 & 0 \\
  0 & 1 & 0 & 0 \\
  0 & 0 & 0 & 1 \\
  0 & 0 & 1 & 0
\end{bmatrix}
\begin{bmatrix}
  1 & 1 & 1 & 1 \\
  1 & -1 & 1 & -1 \\
  1 & 1 & -1 & -1 \\
  1 & -1 & -1 & 1
\end{bmatrix}
=
$$
$$
=
\frac{1}{4}
\begin{bmatrix}
  1 & 1 & 1 & 1 \\
  1 & -1 & -1 & 1 \\
  1 & 1 & -1 & -1 \\
  1 & -1 & 1 & -1
\end{bmatrix}
\begin{bmatrix}
  1 & 1 & 1 & 1 \\
  1 & -1 & 1 & -1 \\
  1 & 1 & -1 & -1 \\
  1 & -1 & -1 & 1
\end{bmatrix}
=
\frac{1}{4}
\begin{bmatrix}
  4 & 0 & 0 & 0 \\
  0 & 0 & 0 & 4 \\
  0 & 0 & 4 & 0 \\
  0 & 4 & 0 & 0
\end{bmatrix} =
\begin{bmatrix}
  1 & 0 & 0 & 0 \\
  0 & 0 & 0 & 1 \\
  0 & 0 & 1 & 0 \\
  0 & 1 & 0 & 0
\end{bmatrix}
$$

\subsubsection{b}

$$
C_{12} \otimes I_3 =
\begin{bmatrix}
  1 I & 0 & 0 & 0 \\
  0 & 1 I & 0 & 0 \\
  0 & 0 & 0 & 1 I \\
  0 & 0 & 1 I & 0
\end{bmatrix} =
\begin{bmatrix}
  1 & 0 & 0 & 0 & 0 & 0 & 0 & 0 \\
  0 & 1 & 0 & 0 & 0 & 0 & 0 & 0 \\
  0 & 0 & 1 & 0 & 0 & 0 & 0 & 0 \\
  0 & 0 & 0 & 1 & 0 & 0 & 0 & 0 \\
  0 & 0 & 0 & 0 & 0 & 0 & 1 & 0 \\
  0 & 0 & 0 & 0 & 0 & 0 & 0 & 1 \\
  0 & 0 & 0 & 0 & 1 & 0 & 0 & 0 \\
  0 & 0 & 0 & 0 & 0 & 1 & 0 & 0 \\
\end{bmatrix}
\quad
C_{13} \otimes I_2 =
\begin{bmatrix}
  1 & 0 & 0 & 0 & 0 & 0 & 0 & 0 \\
  0 & 1 & 0 & 0 & 0 & 0 & 0 & 0 \\
  0 & 0 & 1 & 0 & 0 & 0 & 0 & 0 \\
  0 & 0 & 0 & 1 & 0 & 0 & 0 & 0 \\
  0 & 0 & 0 & 0 & 0 & 1 & 0 & 0 \\
  0 & 0 & 0 & 0 & 1 & 0 & 0 & 0 \\
  0 & 0 & 0 & 0 & 0 & 0 & 0 & 1 \\
  0 & 0 & 0 & 0 & 0 & 0 & 1 & 0 \\
\end{bmatrix}
$$
$$
H_1 \otimes I_2 \otimes I_3 =
\begin{bmatrix}
  1 I & 1 I \\
  1 I & -1 I
\end{bmatrix}
\otimes I_3 =
\begin{bmatrix}
  1 I & 0 & 1I & 0 \\
  0 & 1 I & 0 & 1I \\
  1 I & 0 & -1I & 0 \\
  0 & 1I & 0 & -1I
\end{bmatrix} =
\begin{bmatrix}
  1 & 0 & 0 & 0 & 1 & 0 & 0 & 0 \\
  0 & 1 & 0 & 0 & 0 & 1 & 0 & 0 \\
  0 & 0 & 1 & 0 & 0 & 0 & 1 & 0 \\
  0 & 0 & 0 & 1 & 0 & 0 & 0 & 1 \\
  1 & 0 & 0 & 0 & 1 & 0 & 0 & 0 \\
  0 & 1 & 0 & 0 & 0 & 1 & 0 & 0 \\
  0 & 0 & 1 & 0 & 0 & 0 & 1 & 0 \\
  0 & 0 & 0 & 1 & 0 & 0 & 0 & 1 \\
\end{bmatrix}
$$

$$
(C_{13} \otimes I_2)(C_{12} \otimes I_3)(H_1 \otimes I_2 \otimes I_3) =
$$
$$
=
\begin{bmatrix}
  1 & 0 & 0 & 0 & 0 & 0 & 0 & 0 \\
  0 & 1 & 0 & 0 & 0 & 0 & 0 & 0 \\
  0 & 0 & 1 & 0 & 0 & 0 & 0 & 0 \\
  0 & 0 & 0 & 1 & 0 & 0 & 0 & 0 \\
  0 & 0 & 0 & 0 & 0 & 0 & 1 & 0 \\
  0 & 0 & 0 & 0 & 0 & 0 & 0 & 1 \\
  0 & 0 & 0 & 0 & 1 & 0 & 0 & 0 \\
  0 & 0 & 0 & 0 & 0 & 1 & 0 & 0 \\
\end{bmatrix}
\begin{bmatrix}
  1 & 0 & 0 & 0 & 0 & 0 & 0 & 0 \\
  0 & 1 & 0 & 0 & 0 & 0 & 0 & 0 \\
  0 & 0 & 1 & 0 & 0 & 0 & 0 & 0 \\
  0 & 0 & 0 & 1 & 0 & 0 & 0 & 0 \\
  0 & 0 & 0 & 0 & 0 & 1 & 0 & 0 \\
  0 & 0 & 0 & 0 & 1 & 0 & 0 & 0 \\
  0 & 0 & 0 & 0 & 0 & 0 & 0 & 1 \\
  0 & 0 & 0 & 0 & 0 & 0 & 1 & 0 \\
\end{bmatrix}
\begin{bmatrix}
  1 & 0 & 0 & 0 & 1 & 0 & 0 & 0 \\
  0 & 1 & 0 & 0 & 0 & 1 & 0 & 0 \\
  0 & 0 & 1 & 0 & 0 & 0 & 1 & 0 \\
  0 & 0 & 0 & 1 & 0 & 0 & 0 & 1 \\
  1 & 0 & 0 & 0 & 1 & 0 & 0 & 0 \\
  0 & 1 & 0 & 0 & 0 & 1 & 0 & 0 \\
  0 & 0 & 1 & 0 & 0 & 0 & 1 & 0 \\
  0 & 0 & 0 & 1 & 0 & 0 & 0 & 1 \\
\end{bmatrix}=
$$
$$
=
\begin{bmatrix}
  1 & 0 & 0 & 0 & 0 & 0 & 0 & 0 \\
  0 & 1 & 0 & 0 & 0 & 0 & 0 & 0 \\
  0 & 0 & 1 & 0 & 0 & 0 & 0 & 0 \\
  0 & 0 & 0 & 1 & 0 & 0 & 0 & 0 \\
  0 & 0 & 0 & 0 & 0 & 0 & 0 & 1 \\
  0 & 0 & 0 & 0 & 0 & 0 & 1 & 0 \\
  0 & 0 & 0 & 0 & 0 & 1 & 0 & 0 \\
  0 & 0 & 0 & 0 & 1 & 0 & 0 & 0 \\
\end{bmatrix}
\begin{bmatrix}
  1 & 0 & 0 & 0 & 1 & 0 & 0 & 0 \\
  0 & 1 & 0 & 0 & 0 & 1 & 0 & 0 \\
  0 & 0 & 1 & 0 & 0 & 0 & 1 & 0 \\
  0 & 0 & 0 & 1 & 0 & 0 & 0 & 1 \\
  1 & 0 & 0 & 0 & 1 & 0 & 0 & 0 \\
  0 & 1 & 0 & 0 & 0 & 1 & 0 & 0 \\
  0 & 0 & 1 & 0 & 0 & 0 & 1 & 0 \\
  0 & 0 & 0 & 1 & 0 & 0 & 0 & 1 \\
\end{bmatrix}=
\begin{bmatrix}
  1 & 0 & 0 & 0 & 0 & 0 & 0 & 0 \\
  0 & 1 & 0 & 0 & 0 & 0 & 0 & 0 \\
  0 & 0 & 1 & 0 & 0 & 0 & 0 & 0 \\
  0 & 0 & 0 & 1 & 0 & 0 & 0 & 0 \\
  0 & 0 & 0 & 1 & 0 & 0 & 0 & 1 \\
  0 & 0 & 1 & 0 & 0 & 0 & 1 & 0 \\
  0 & 1 & 0 & 0 & 0 & 1 & 0 & 0 \\
  1 & 0 & 0 & 0 & 1 & 0 & 0 & 0 \\
\end{bmatrix}
$$

\subsubsection{c}

$$
S_{12} \otimes I_3 =
\begin{bmatrix}
  1 I & 0 & 0 & 0 \\
  0 & 0 & 1 I & 0 \\
  0 & 1 I & 0 & 0 \\
  0 & 0 & 0 & 1 I
\end{bmatrix} =
\begin{bmatrix}
  1 & 0 & 0 & 0 & 0 & 0 & 0 & 0 \\
  0 & 1 & 0 & 0 & 0 & 0 & 0 & 0 \\
  0 & 0 & 0 & 0 & 1 & 0 & 0 & 0 \\
  0 & 0 & 0 & 0 & 0 & 1 & 0 & 0 \\
  0 & 0 & 1 & 0 & 0 & 0 & 0 & 0 \\
  0 & 0 & 0 & 1 & 0 & 0 & 0 & 0 \\
  0 & 0 & 0 & 0 & 0 & 0 & 1 & 0 \\
  0 & 0 & 0 & 0 & 0 & 0 & 0 & 1 \\
\end{bmatrix}
$$
$$
H_1 \otimes I_2 \otimes H_3 =
\begin{bmatrix}
  1 & 0 & 1 & 0 \\
  0 & 1 & 0 & 1 \\
  1 & 0 & -1 & 0 \\
  0 & 1 & 0 & -1
\end{bmatrix} \otimes H_3 =
\begin{bmatrix}
  1 & 1 & 0 & 0 & 1 & 1 & 0 & 0 \\
  1 & -1 & 0 & 0 & 1 & -1 & 0 & 0 \\
  0 & 0 & 1 & 1 & 0 & 0 & 1 & 1 \\
  0 & 0 & 1 & -1 & 0 & 0 & 1 & -1 \\
  1 & 1 & 0 & 0 & -1 & -1 & 0 & 0 \\
  1 & -1 & 0 & 0 & -1 & 1 & 0 & 0 \\
  0 & 0 & 1 & 1 & 0 & 0 & 1 & 1 \\
  0 & 0 & 1 & -1 & 0 & 0 & 1 & -1 \\
\end{bmatrix}
$$

$$
(C_{13} \otimes I_2)(S_{12} \otimes I_3)(H_1 \otimes I_2 \otimes H_3) =
$$
$$
=
\begin{bmatrix}
  1 & 0 & 0 & 0 & 0 & 0 & 0 & 0 \\
  0 & 1 & 0 & 0 & 0 & 0 & 0 & 0 \\
  0 & 0 & 1 & 0 & 0 & 0 & 0 & 0 \\
  0 & 0 & 0 & 1 & 0 & 0 & 0 & 0 \\
  0 & 0 & 0 & 0 & 0 & 1 & 0 & 0 \\
  0 & 0 & 0 & 0 & 1 & 0 & 0 & 0 \\
  0 & 0 & 0 & 0 & 0 & 0 & 0 & 1 \\
  0 & 0 & 0 & 0 & 0 & 0 & 1 & 0 \\
\end{bmatrix}
\begin{bmatrix}
  1 & 0 & 0 & 0 & 0 & 0 & 0 & 0 \\
  0 & 1 & 0 & 0 & 0 & 0 & 0 & 0 \\
  0 & 0 & 0 & 0 & 1 & 0 & 0 & 0 \\
  0 & 0 & 0 & 0 & 0 & 1 & 0 & 0 \\
  0 & 0 & 1 & 0 & 0 & 0 & 0 & 0 \\
  0 & 0 & 0 & 1 & 0 & 0 & 0 & 0 \\
  0 & 0 & 0 & 0 & 0 & 0 & 1 & 0 \\
  0 & 0 & 0 & 0 & 0 & 0 & 0 & 1 \\
\end{bmatrix}
\begin{bmatrix}
  1 & 1 & 0 & 0 & 1 & 1 & 0 & 0 \\
  1 & -1 & 0 & 0 & 1 & -1 & 0 & 0 \\
  0 & 0 & 1 & 1 & 0 & 0 & 1 & 1 \\
  0 & 0 & 1 & -1 & 0 & 0 & 1 & -1 \\
  1 & 1 & 0 & 0 & -1 & -1 & 0 & 0 \\
  1 & -1 & 0 & 0 & -1 & 1 & 0 & 0 \\
  0 & 0 & 1 & 1 & 0 & 0 & 1 & 1 \\
  0 & 0 & 1 & -1 & 0 & 0 & 1 & -1 \\
\end{bmatrix}=
$$
$$
=
\begin{bmatrix}
  1 & 0 & 0 & 0 & 0 & 0 & 0 & 0 \\
  0 & 1 & 0 & 0 & 0 & 0 & 0 & 0 \\
  0 & 0 & 0 & 0 & 1 & 0 & 0 & 0 \\
  0 & 0 & 0 & 0 & 0 & 1 & 0 & 0 \\
  0 & 0 & 0 & 1 & 0 & 0 & 0 & 0 \\
  0 & 0 & 1 & 0 & 0 & 0 & 0 & 0 \\
  0 & 0 & 0 & 0 & 0 & 0 & 0 & 1 \\
  0 & 0 & 0 & 0 & 0 & 0 & 1 & 0 \\
\end{bmatrix}
\begin{bmatrix}
  1 & 1 & 0 & 0 & 1 & 1 & 0 & 0 \\
  1 & -1 & 0 & 0 & 1 & -1 & 0 & 0 \\
  0 & 0 & 1 & 1 & 0 & 0 & 1 & 1 \\
  0 & 0 & 1 & -1 & 0 & 0 & 1 & -1 \\
  1 & 1 & 0 & 0 & -1 & -1 & 0 & 0 \\
  1 & -1 & 0 & 0 & -1 & 1 & 0 & 0 \\
  0 & 0 & 1 & 1 & 0 & 0 & 1 & 1 \\
  0 & 0 & 1 & -1 & 0 & 0 & 1 & -1 \\
\end{bmatrix}
=
\begin{bmatrix}
  1 & 1 & 0 & 0 & 1 & 1 & 0 & 0 \\
  1 & -1 & 0 & 0 & 1 & -1 & 0 & 0 \\
  1 & 1 & 0 & 0 & -1 & -1 & 0 & 0 \\
  1 & -1 & 0 & 0 & -1 & 1 & 0 & 0 \\
  0 & 0 & 1 & -1 & 0 & 0 & 1 & -1 \\
  0 & 0 & 1 & 1 & 0 & 0 & 1 & 1 \\
  0 & 0 & 1 & -1 & 0 & 0 & 1 & -1 \\
  0 & 0 & 1 & 1 & 0 & 0 & 1 & 1 \\
\end{bmatrix}
$$

\section{Zestaw 3}

\subsection{Zadanie 1}

Z definicji normy Frobeniusa:
$$
||A||_F = \sqrt{\sum_{i=1}^{m}\sum_{j=1}^{n} |a_{ij}|^2}
$$
Mamy udowodnić:
$$
||A||_F^2 = \Tr(A^\dag A) = \sum_i \sigma_i^2
$$

$$
A^\dag A =
\begin{bmatrix}
  a^*_{11} & a^*_{21} & \dots & a^*_{1n} \\
  a^*_{12} & a^*_{22} & \dots & a^*_{2n} \\
  \vdots & \vdots & \ddots & \vdots \\
  a^*_{m1} & a^*_{m2} & \dots & a^*_{mn} \\
\end{bmatrix}
\begin{bmatrix}
  a_{11} & a_{12} & \dots & a_{1m} \\
  a_{21} & a_{22} & \dots & a_{2m} \\
  \vdots & \vdots & \ddots & \vdots \\
  a_{n1} & a_{n2} & \dots & a_{nm} \\
\end{bmatrix}
\Rightarrow
(A^\dag A)_{jj} = \sum_{i=1}^{m} a^*_{ij} a_{ij} = \sum_{i=1}^{m} |a_{ij}|^2
$$
$$
\Tr(A^\dag A) = \sum_{j=1}^{n} (A^\dag A)_{jj} = \sum_{j=1}^{n}
\sum_{i=1}^{m} |a_{ij}|^2 = ||A||_F^2
$$

Z definicji SVD, mamy:
$$
A = U \Sigma V^\dag
$$
gdzie $U$ i $V$ są ortogonalnymi macierzami unitarnymi, a $\Sigma$
jest diagonalną macierzą wartości własnych.
$$
(A^\dag A) = (U \Sigma V^\dag)^\dag (U \Sigma V^\dag) = V \Sigma^\dag
U^\dag U \Sigma V^\dag
$$
z ortogonalności $U$:
$$
(A^\dag A) = V \Sigma^\dag I \Sigma V^\dag = V \Sigma^\dag \Sigma V^\dag
$$
Ponieważ $\Tr$ jest niezmienniczy względem kolejności operacji:
$$
\Tr(A^\dag A) = \Tr(V \Sigma^\dag \Sigma V^\dag) = \Tr(\Sigma^\dag
\Sigma V^\dag V) = \Tr(\Sigma^\dag \Sigma I)
$$
A ponieważ, $\Sigma$ jest diagonalną macierzą wartości osobliwych, to mamy:
$$
\Tr(\Sigma^\dag \Sigma) = \sum_{i=1}^{n} \sigma_i^2 = ||A||_F^2 \quad \square
$$

\subsection{Zadanie 2}
$$
\Tr_B(\ket{\Psi}\bra{\Phi} \otimes \ket{\alpha}\bra{\beta}) =
\braket{\beta}{\alpha} \ket{\Psi}\bra{\Phi}
$$

$$
\Tr_B(\rho_{AB}) = \sum_k (I_A \otimes \bra{k}_B) \rho_{AB} (I_A
\otimes \ket{k}_B)
$$
$$
\Tr_B(\ket{\Psi}\bra{\Phi} \otimes \ket{\alpha}\bra{\beta}) = \sum_k
(I_A \otimes \bra{k}_B)\left(\ket{\Psi}\bra{\Phi} \otimes
\ket{\alpha}\bra{\beta}\right) (I_A
\otimes \ket{k}_B) =
$$
$$
=\sum_k \left[(I_A \ket{\Psi}\bra{\Phi}) \otimes (\bra{k}_B
\ket{\alpha}\bra{\beta})\right] (I_A
\otimes \ket{k}_B) = \sum_k (\ket{\Psi}\bra{\Phi} I_A) \otimes (\bra{k}_B
\ket{\alpha}\bra{\beta} \ket{k}_B) = \sum_k \ket{\Psi}\bra{\Phi}
\otimes (\bra{k}_B
\ket{\alpha}\bra{\beta} \ket{k}_B)
$$
Ponieważ $\braket{\cdot}{\cdot}$ to skalar, $\otimes$ redukuje się do mnożenia,
więc:
$$
\Tr_B(\ket{\Psi}\bra{\Phi} \otimes \ket{\alpha}\bra{\beta}) =
\ket{\Psi}\bra{\Phi} \sum_k \bra{k}_B \ket{\alpha}\bra{\beta}
\ket{k}_B = \ket{\Psi}\bra{\Phi} \sum_k \braket{\alpha}{k} \braket{k}{\beta}
$$
Relacja zupełności bazy ortonormalnej mówi nam, że:
$$
\sum_k \ket{k}\bra{k} = I \Rightarrow \Tr_B(\ket{\Psi}\bra{\Phi}
\otimes \ket{\alpha}\bra{\beta}) = \ket{\Psi}\bra{\Phi}
\braket{\alpha}{\beta} \quad \square
$$

\subsection{Zadanie 3}

$$
U(a) = e^{iA}, A = A^\dag \rightarrow U(A)^\dag U(A) = I =
U(A)U(A)^\dag
$$

Z definicji:
$$
U = e^{iA} = \sum_{n=1}^{\infty} \frac{(iA)^n}{n!} \Rightarrow U^\dag
= \sum_{n=1}^{\infty} \left(\frac{(iA)^n}{n!}\right)^\dag =
\sum_{n=1}^{\infty} \frac{((iA)^n)^\dag}{n!}
$$
Ponieważ $A$ jest hermitowska:
$$
((iA)^n)^\dag = (A^\dag)^n(i^n)^\dag = A^n (-i)^n = (-iA)^n
$$
$$
U^\dag = \sum_{n=1}^{\infty} \frac{(-iA)^n}{n!} = e^{-iA}
$$
$$
U^\dag U = e^{-iA} e^{iA} = I \quad UU^\dag = e^{iA} e^{-iA} = I \quad \square
$$

\subsection{Zadanie 4}

$$
A = A^\dag \quad A\ket{\alpha} = \alpha \ket{\alpha} \quad A\ket{\beta} = \beta
\ket{\beta} \quad \braket{\alpha}{\beta} = 0 \quad \alpha \neq \beta
$$

$$
\bra{\beta}A\ket{\alpha} = \alpha \braket{\beta}{\alpha}
$$
$$
\bra{\beta}A\ket{\alpha} = (A^\dag \ket{\beta})^\dag \ket{\alpha} =
\beta \braket{\beta}{\alpha}
$$
$$
\alpha \braket{\beta}{\alpha} = \beta \braket{\beta}{\alpha}
$$
Więc, ponieważ $\alpha \neq \beta$, $\braket{\alpha}{\beta} = 0 \quad \square$

\end{document}
